% GENERATED FILE. Edit canonical lessons, then run npm run generate:books.
% canonical-source-hash: fnv1a64:27b7998c8eab4dab
% canonical-lessons: HI-C344-parinam, HI-C344-prabhav, HI-C344-barhna, HI-C344-bigaarna, HI-C344-srot

\chapter{What Follows From It}
\label{ch:hi-what-follows-from-it}
\hlchaptermodality{344}

\begin{chapteropening}
\textbf{By the end of this chapter:} \emph{I can say what an event led to.}

Complete the last lesson of chapter 344: I can say what an event led to.
\end{chapteropening}

\section[\texorpdfstring{\hi{परिणाम} (pariṇām)}{परिणाम (pariṇām)}]{\hi{परिणाम} (pariṇām) — a consequence, an outcome}

\label{lesson:HI-C344-parinam}



\begin{quote}
Before the new one: say the Hindi for a basis, grounds, then the Hindi for responsible.
\end{quote}

\subsection*{You'll want to know: \hi{परिणाम}}
\textbf{\hi{परिणाम}} — \emph{pariṇām} — \textquotedblleft{}a consequence, an outcome\textquotedblright{}.

\begin{quote}
\hi{हर} \hi{काम} \hi{का} \hi{कोई} \hi{परिणाम} \hi{होता} \hi{है।}
\end{quote}

\begin{quote}
\emph{har kām kā koī pariṇām hotā hai.}
\end{quote}

\begin{quote}
Every action has some consequence.
\end{quote}

Masculine.

\begin{cousinweb}
\textbf{\hi{परिणाम}} (\emph{pariṇām}) is a tatsama, borrowed from Sanskrit \emph{pariṇāma}, \textquotedblleft{}a change; a result\textquotedblright{} (Monier-Williams).
\end{cousinweb}

\begin{grammarlens}[title={What you've built}]
One more word for giving reasons and opinions.
\end{grammarlens}

\subsection*{Your turn}
\begin{itemize}
\raggedright
  \item \emph{Say it:} \emph{pariṇām}
  \item \emph{Say it:} \emph{pariṇām}, once more
  \item \emph{Say it:} \emph{har kām kā koī pariṇām hotā hai.} — Every action has some consequence.
  \item \emph{Recall:} say the Hindi for a basis, grounds, then the Hindi for responsible, then say \emph{pariṇām} again
\end{itemize}

\begin{tcolorbox}[breakable,colback=teal!4,colframe=teal!35!black,title={Before you move on}]
What does \hi{परिणाम} mean? (\textquotedblleft{}A consequence, an outcome\textquotedblright{}.) Say it once more.
\end{tcolorbox}
\section[\texorpdfstring{\hi{प्रभाव} (prabhāv)}{प्रभाव (prabhāv)}]{\hi{प्रभाव} (prabhāv) — an influence}

\label{lesson:HI-C344-prabhav}



\begin{quote}
Before the new one: say the Hindi for responsible, then the Hindi for a consequence.
\end{quote}

\subsection*{You'll want to know: \hi{प्रभाव}}
\textbf{\hi{प्रभाव}} — \emph{prabhāv} — \textquotedblleft{}an influence\textquotedblright{}.

\begin{quote}
\hi{बच्चों} \hi{पर} \hi{दोस्तों} \hi{का} \hi{बहुत} \hi{प्रभाव} \hi{होता} \hi{है।}
\end{quote}

\begin{quote}
\emph{baccoṁ par dostoṁ kā bahut prabhāv hotā hai.}
\end{quote}

\begin{quote}
Friends have a great influence on children.
\end{quote}

Masculine.

\begin{cousinweb}
\textbf{\hi{प्रभाव}} (\emph{prabhāv}) is a tatsama, borrowed from Sanskrit \emph{prabhāva}, \textquotedblleft{}might, power\textquotedblright{} (Monier-Williams).
\end{cousinweb}

\begin{grammarlens}[title={What you've built}]
One more word for giving reasons and opinions.
\end{grammarlens}

\subsection*{Your turn}
\begin{itemize}
\raggedright
  \item \emph{Say it:} \emph{prabhāv}
  \item \emph{Say it:} \emph{prabhāv}, once more
  \item \emph{Say it:} \emph{baccoṁ par dostoṁ kā bahut prabhāv hotā hai.} — Friends have a great influence on children.
  \item \emph{Recall:} say the Hindi for responsible, then the Hindi for a consequence, then say \emph{prabhāv} again
\end{itemize}

\begin{tcolorbox}[breakable,colback=teal!4,colframe=teal!35!black,title={Before you move on}]
What does \hi{प्रभाव} mean? (\textquotedblleft{}An influence\textquotedblright{}.) Say it once more.
\end{tcolorbox}
\section[\texorpdfstring{\hi{बढ़ना} (baṛhnā)}{बढ़ना (baṛhnā)}]{\hi{बढ़ना} (baṛhnā) — to rise, to go up (of numbers, prices)}

\label{lesson:HI-C344-barhna}



\begin{quote}
Before the new one: say the Hindi for a consequence, then the Hindi for an influence.
\end{quote}

\subsection*{You'll want to know: \hi{बढ़ना}}
\textbf{\hi{बढ़ना}} — \emph{baṛhnā} — \textquotedblleft{}to rise, to go up (of numbers, prices)\textquotedblright{}.

\begin{quote}
\hi{गर्मी} \hi{में} \hi{बिजली} \hi{का} \hi{बिल} \hi{बढ़ता} \hi{है।}
\end{quote}

\begin{quote}
\emph{garmī meṁ bijlī kā bil baṛhtā hai.}
\end{quote}

\begin{quote}
The electricity bill goes up in summer.
\end{quote}

Past (it …): \hi{वह} \hi{किराया} \hi{बढ़} \hi{गया}, \emph{vah kirāyā baṛh gayā}.

\begin{cousinweb}
\textbf{\hi{बढ़ना}} (\emph{baṛhnā}) is a tadbhava, inherited from Sanskrit \emph{vardhate}, \textquotedblleft{}grows, increases\textquotedblright{} (CDIAL 11376).
\end{cousinweb}

\begin{grammarlens}[title={What you've built}]
One more word for giving reasons and opinions.
\end{grammarlens}

\subsection*{Your turn}
\begin{itemize}
\raggedright
  \item \emph{Say it:} \emph{baṛhnā}
  \item \emph{Say it:} \emph{baṛhnā}, once more
  \item \emph{Say it:} \emph{garmī meṁ bijlī kā bil baṛhtā hai.} — The electricity bill goes up in summer.
  \item \emph{Recall:} say the Hindi for a consequence, then the Hindi for an influence, then say \emph{baṛhnā} again
\end{itemize}

\begin{tcolorbox}[breakable,colback=teal!4,colframe=teal!35!black,title={Before you move on}]
What does \hi{बढ़ना} mean? (\textquotedblleft{}To rise, to go up (of numbers, prices)\textquotedblright{}.) Say it once more.
\end{tcolorbox}
\section[\texorpdfstring{\hi{बिगाड़ना} (bigāṛnā)}{बिगाड़ना (bigāṛnā)}]{\hi{बिगाड़ना} (bigāṛnā) — to spoil, to make worse}

\label{lesson:HI-C344-bigaarna}



\begin{quote}
Before the new one: say the Hindi for an influence, then the Hindi for to rise.
\end{quote}

\subsection*{You'll want to know: \hi{बिगाड़ना}}
\textbf{\hi{बिगाड़ना}} — \emph{bigāṛnā} — \textquotedblleft{}to spoil, to make worse\textquotedblright{}.

\begin{quote}
\hi{बारिश} \hi{ने} \hi{फ़सल} \hi{बिगाड़} \hi{दी।}
\end{quote}

\begin{quote}
\emph{bāriś ne fasal bigāṛ dī.}
\end{quote}

\begin{quote}
The rain ruined the crop.
\end{quote}

Past, with \hi{ने} (he or she …): \hi{उसने} \hi{खेल} \hi{बिगाड़} \hi{दिया}, \emph{usne khel bigāṛ diyā}.

\begin{cousinweb}
\textbf{\hi{बिगाड़ना}} (\emph{bigāṛnā}) is a tadbhava, inherited from Sanskrit \emph{vighaṭate}, \textquotedblleft{}breaks apart\textquotedblright{} (CDIAL 11673).
\end{cousinweb}

\begin{grammarlens}[title={What you've built}]
One more word for giving reasons and opinions.
\end{grammarlens}

\subsection*{Your turn}
\begin{itemize}
\raggedright
  \item \emph{Say it:} \emph{bigāṛnā}
  \item \emph{Say it:} \emph{bigāṛnā}, once more
  \item \emph{Say it:} \emph{bāriś ne fasal bigāṛ dī.} — The rain ruined the crop.
  \item \emph{Recall:} say the Hindi for an influence, then the Hindi for to rise, then say \emph{bigāṛnā} again
\end{itemize}

\begin{tcolorbox}[breakable,colback=teal!4,colframe=teal!35!black,title={Before you move on}]
What does \hi{बिगाड़ना} mean? (\textquotedblleft{}To spoil, to make worse\textquotedblright{}.) Say it once more.
\end{tcolorbox}
\section[\texorpdfstring{\hi{स्रोत} (srot)}{स्रोत (srot)}]{\hi{स्रोत} (srot) — a source}

\label{lesson:HI-C344-srot}



\begin{quote}
Before the new one: say the Hindi for to rise, then the Hindi for to make worse.
\end{quote}

\subsection*{You'll want to know: \hi{स्रोत}}
\textbf{\hi{स्रोत}} — \emph{srot} — \textquotedblleft{}a source\textquotedblright{}.

\begin{quote}
\hi{ख़बर} \hi{का} \hi{स्रोत} \hi{क्या} \hi{है}?
\end{quote}

\begin{quote}
\emph{khabar kā srot kyā hai?}
\end{quote}

\begin{quote}
What is the source of the news?
\end{quote}

Masculine.

\begin{cousinweb}
\textbf{\hi{स्रोत}} (\emph{srot}) is a tatsama, borrowed from Sanskrit \emph{srotas}, \textquotedblleft{}a stream, a current\textquotedblright{} (Monier-Williams). Hindi uses it for the source of anything, news included.
\end{cousinweb}

\begin{grammarlens}[title={What you've built}]
One more word for giving reasons and opinions.
\end{grammarlens}

\subsection*{Your turn}
\begin{itemize}
\raggedright
  \item \emph{Say it:} \emph{srot}
  \item \emph{Say it:} \emph{srot}, once more
  \item \emph{Say it:} \emph{khabar kā srot kyā hai?} — What is the source of the news?
  \item \emph{Recall:} say the Hindi for to rise, then the Hindi for to make worse, then say \emph{srot} again
\end{itemize}

\begin{tcolorbox}[breakable,colback=teal!4,colframe=teal!35!black,title={Before you move on}]
What does \hi{स्रोत} mean? (\textquotedblleft{}A source\textquotedblright{}.) Say it once more.
\end{tcolorbox}
